\documentclass[10pt,a4paper]{amsart}
\usepackage[margin=19mm]{geometry}
\usepackage{amsmath,amssymb,mathtools}
\usepackage[scr=rsfs,cal=euler]{mathalfa}
\usepackage{xfrac}
\usepackage[unicode,hidelinks,bookmarks=false]{hyperref}
\newcommand{\ie}{i.e.\ }
\newcommand{\cf}{cf.\ }
\newcommand{\resp}{respectively\ }
\newcommand{\Subsection}{\subsection}
\providecommand{\nobreakdash}{\nobreak\hbox{-}}
\setlength{\emergencystretch}{2em}
\multlinegap0pt
\usepackage{bm}
\usepackage{bbm}
\usepackage{dsfont}
\usepackage{xspace}
\usepackage{cancel}
\usepackage{mathtools}
\usepackage{amsthm}
\makeatletter
\@ifundefined{theorem}{\newtheorem{theorem}{Theorem}}{}
\@ifundefined{lemma}{\newtheorem{lemma}[theorem]{Lemma}}{}
\@ifundefined{proposition}{\newtheorem{proposition}[theorem]{Proposition}}{}
\@ifundefined{remark}{\newtheorem{remark}[theorem]{Remark}}{}
\makeatother
\newcommand {\F}{\mathcal{F}}
\newcommand {\Ma}{\mathcal{M}}
\newcommand {\Rn}{\mathbb{R}^{n}}
\newcommand {\Sp}{S^{*}\Rn}
\newcommand {\Sw}{\mathcal{S}}
\newcommand {\W}{\mathrm{W}}
\newcommand {\Z}{\mathbb Z}
\newcommand {\la}{\lambda}
\newcommand {\lb}{\langle}
\newcommand {\ph}{\varphi}
\newcommand {\rb}{\rangle}
\newcommand {\supp}{\mathrm{supp}}
\newcommand {\ud}{\mathrm{d}}
\newcommand {\veps}{\varepsilon}
\newcommand {\wt}{\widetilde}
\newcommand {\w}{\omega}
\title[The Hardy spaces $\mathcal{H}^{p}\_{FIO}(\mathbb{R}^{n})$ for]{The Hardy spaces $\mathcal{H}^{p}\_{FIO}(\mathbb{R}^{n})$ for Fourier integral operators for $p<1$}
\author{Naijia Liu, Jan Rozendaal, Liang Song}
\date{Source: arXiv:2508.13243v1 (CC BY 4.0)}
\begin{document}
\maketitle

\noindent\emph{Source and licence.} The Hardy spaces $\mathcal{H}^{p}_{FIO}(\mathbb{R}^{n})$ for Fourier integral operators for $p<1$. Version arXiv:2508.13243v1 (CC BY 4.0), distributed under the Creative Commons Attribution 4.0 International licence, as recorded in the arXiv licence metadata for this exact version and shown on the abstract page https://arxiv.org/abs/2508.13243v1. Full rights notice: licenses/arxiv-2508.13243-CC-BY-4.0.txt.\par\smallskip

\noindent\textit{Editorial scope.} Counted unit: the Proposition whose source label is \texttt{prop:maxineq} in the article (source0.tex lines 1435--1441, source label \texttt{prop:maxineq}), delivered together with the article's own complete original proof of it (source0.tex lines 1442--1551, one \texttt{proof} environment of 110 lines). No mathematics is written, repaired or re-derived: statement and proof are the article's own text line for line. Required context retained uncounted: eq:equivmetric (lines 434--436); eq:HLcontrol (lines 480--484); eq:phiproperties3 (lines 793--795); lem:packetbounds (lines 805--816); eq:boundspsi (lines 808--810); eq:boundspsiinverse (lines 812--814); rem:thetatilde (lines 818--825). Every definition, display and label of those lines is retained unchanged. Omitted from this package and not redistributed: 1--433, 437--479, 485--792, 796--804, 811--811, 815--817, 826--1434, 1552--1855. Nothing inside a retained excerpt is omitted or altered: every retained body is delivered exactly as the article's lines are written, so no omission receipt of any kind accompanies this package. Citations: the retained excerpts carry no citation command at all, so no bibliography entry of the article is transcribed and no reference is redistributed. Reference adaptation: none was needed and none was made; every reference inside the delivered unit resolves inside it, so no unresolved reference or citation is produced. Printed slips kept literal: no printed word, symbol, hypothesis or number of the counted statement or of its proof was altered. Layout and class: the article's own class is \texttt{amsart} with packages including ae, aecompl, esint, graphicx, xy, amsmath, amscd, booktabs, amsthm, amssymb, amsfonts, qsymbols, latexsym, mathrsfs; the wrapper is the lane's pinned \texttt{amsart} editorial wrapper at 10pt on A4, which restates the theorem-like environments the retained text needs and adds the article's own macro definitions that the retained text needs (\texttt{\textbackslash F}, \texttt{\textbackslash Ma}, \texttt{\textbackslash Rn}, \texttt{\textbackslash Sp}, \texttt{\textbackslash Sw}, \texttt{\textbackslash W}, \texttt{\textbackslash Z}, \texttt{\textbackslash la}, \texttt{\textbackslash lb}, \texttt{\textbackslash ph}, \texttt{\textbackslash rb}, \texttt{\textbackslash supp}, \texttt{\textbackslash ud}, \texttt{\textbackslash veps}, \texttt{\textbackslash w}, \texttt{\textbackslash wt}), with extra wrapper packages (bm, bbm, dsfont, xspace, cancel, mathtools). The delivered numbering is the unit's own. Assets: no retained span contains a figure environment, a graphics-inclusion command, a PDF-page inclusion, a table or a TikZ picture; no image file is copied into this package, the package declares no asset dependency and no image file is redistributed with it. Licence: version arXiv:2508.13243v1 of this article is distributed under the Creative Commons Attribution 4.0 International licence, as recorded in the arXiv licence metadata for this exact version and shown on the abstract page https://arxiv.org/abs/2508.13243v1. A full rights notice is delivered at licenses/arxiv-2508.13243-CC-BY-4.0.txt. Characters: every retained excerpt is pure ASCII, so the delivered engine, XeLaTeX with its default Latin Modern text font, reports no missing character.
\begin{equation}\label{eq:equivmetric}
d((x,\omega),(y,\nu))\eqsim \big{(}|\omega\cdot(x-y)|+|x-y|^{2}+|\omega-\nu|^{2}\big{)}^{1/2}
\end{equation}

\begin{align}\label{eq:HLcontrol}
{\sigma}^{-n}\int_{\Sp}(1+{\sigma}^{-1}d((x,\omega),(y,\nu))^{2})^{-M}
|g(y,\nu)|\ud y\ud\nu
\leq C_{M}\Ma(g)(x,\omega)
\end{align}

\begin{equation}\label{eq:phiproperties3}
f=\int_{S^{n-1}}m(D)\ph_{\nu}(D)f\ud\nu.
\end{equation}

\begin{lemma}\label{lem:packetbounds}
For $\w\in S^{n-1}$ and $\sigma\in(0,1)$, let $\eta_{\w,\sigma}\in\{\psi_{\w,\sigma},\theta_{\w,\sigma}\}$. Then $\eta_{\w,\sigma}\in C^{\infty}_{c}(\Rn)$, and each $\xi\in\supp(\eta_{\w,\sigma})$ satisfies $\frac{1}{2}\sigma^{-1}\leq |\xi|\leq 2\sigma^{-1}$ and $|\hat{\xi}-\w|\leq 2\sqrt{\sigma}$. 
Moreover, for all $\alpha\in\Z_{+}^{n}$ and $\beta\in\Z_{+}$ there exists a $C_{\alpha,\beta}\geq0$, independent of $\w$ and $\sigma$, such that
\begin{equation}\label{eq:boundspsi}
|(\w\cdot\partial_{\xi})^{\beta}\partial_{\xi}^{\alpha}\eta_{\w,\sigma}(\xi)|\leq C_{\alpha,\beta}\sigma^{-\frac{n-1}{4}+\frac{|\alpha|}{2}+\beta}
\end{equation}
for all $\xi\in\Rn$. Also, for each $N\geq0$ there exists a $C_{N}\geq0$, independent of $\w$ and $\sigma$, such that
\begin{equation}\label{eq:boundspsiinverse}
|\F^{-1}(\eta_{\w,\sigma})(x)|\leq C_{N}\sigma^{-\frac{3n+1}{4}}(1+\sigma^{-1}|x|^{2}+\sigma^{-2}(\w\cdot x)^{2})^{-N}
\end{equation}
for all $x\in\Rn$.
\end{lemma}

\begin{remark}\label{rem:thetatilde}
Let $\wt{\Psi}\in C^{\infty}_{c}(\Rn)$ be such that $\wt{\Psi}(\xi)=0$ for $|\xi|\notin [\frac{1}{4},4]$, and $\wt{\Psi}\equiv 1$ on $\supp(\Psi)$. For $\w\in S^{n-1}$, $\sigma>0$ and $\xi\in\Rn$, set
\[
\tilde{\theta}_{\w,\sigma}(\xi):=\sigma^{\frac{n-1}{4}}m(\xi)\ph_{\w}(\xi)\wt{\Psi}(\sigma\xi),
\]
where $m$ is as in \eqref{eq:phiproperties3}.
Then the $\tilde{\theta}_{\w,\sigma}$ have the same properties as in Lemma \ref{lem:packetbounds}, with the exception that each $\xi\in\supp(\tilde{\theta}_{\w,\sigma})$ satisfies $\frac{1}{4}\sigma^{-1}\leq |\xi|\leq 4\sigma^{-1}$. More generally, \eqref{eq:boundspsi} and \eqref{eq:boundspsiinverse} still hold if one imposes slightly different support conditions on the $\eta_{\w,\sigma}$.
\end{remark}

\begin{proposition}\label{prop:maxineq}
Let $\lambda>0$ and $N>n/\la$. Then there exists a $C\geq0$ such that
\[
M_{N,\sigma}^{*}(W_{\sigma}f)(x,\omega)\leq C \mathcal{M}_{\lambda}(W_{\sigma}f)(x,\omega)
\]
for all $f\in \Sw'(\Rn)$, $(x,\w)\in\Sp$ and $\sigma\in(0,1)$.
\end{proposition}

\begin{proof}
Fix $f\in\Sw'(\Rn)$, $(x,\w)\in\Sp$ and $\sigma\in (0,1)$. %It is worth pointing out that one could replace $W_{\sigma}f$ in the statement by a general $g_{\sigma}:\Sp\to\C$ satisfying 
%\begin{equation}\label{eq:supportgsigma}
%\supp(\F g_{\sigma}(\cdot,\w))\subseteq\{\xi\in\Rn\mid \tfrac{1}{2}\sigma^{-1}\leq |\xi|\leq 2\sigma^{-1},|\hat{\xi}-\w|\leq 2\sqrt{\sigma}\}.
%\end{equation}
%After a minor modification, this support condition is all that is needed for the proof.
%
%Now, f
First note that, by \eqref{eq:phiproperties3} and Remark \ref{rem:thetatilde}, 
\begin{equation}\label{eq:reprotheta}
%\begin{align}\label{eqn-21-10-6-12}
\theta_{\nu,\sigma}(D)f(y)=\sigma^{-\frac{n-1}{4}}\int_{S^{n-1}}\tilde{\theta}_{\nu,\sigma}(D)\theta_{\mu,\sigma}(D)f(y) \ud\mu
%\end{align}
\end{equation}
for all $(y,\nu)\in\Sp$. In fact, by Lemma \ref{lem:packetbounds} and Remark \ref{rem:thetatilde}, the integrand is only nonzero where $|\mu-\nu|\leq 4\sqrt{\sigma}$. In the latter case, Remark \ref{rem:thetatilde} and \eqref{eq:equivmetric} yield
\begin{align*}
|\F^{-1}(\tilde{\theta}_{\nu,\sigma})(y-z)|&\lesssim \sigma^{-\frac{3n+1}{4}}(1+\sigma^{-1}d((y,\nu),(z,\nu))^{2})^{-N}\\
&\lesssim \sigma^{-\frac{3n+1}{4}}(1+\sigma^{-1}d((y,\nu),(z,\mu))^{2})^{-N}
\end{align*}
for all $z\in\Rn$. Hence we can use the triangle inequality to write
\begin{align*}
&(1+\sigma^{-1}d((x,\omega),(y,\nu))^{2})^{-N}|\theta_{\nu,\sigma}(D)f(y)|\\
&\lesssim\sigma^{-n}\int_{\Sp}\frac{|\theta_{\mu,{\sigma}}(D)f(z)|}{(1+\sigma^{-1}d((x,\omega),(y,\nu))^{2})^{N}(1+\sigma^{-1}d((y,\nu),(z,\mu))^{2})^{N}}
\ud z\ud\mu\\
&\lesssim \sigma^{-n}\int_{\Sp}\frac{|
\theta_{\mu,\sigma}(D)f(z)|^{1-\lambda}|\theta_{\mu,\sigma}(D)f(z)|^{\lambda} }{(1+\sigma^{-1}d((x,\omega),(z,\mu))^{2})^{N}}\ud z\ud\mu.
\end{align*}
By combining this with \eqref{eq:HLcontrol}, we see that
\begin{align*}
&M_{N,\sigma}^{*}(W_{\sigma}f)(x,\omega)\\
&\lesssim \sigma^{-n}\int_{\Sp}\frac{|
\theta_{\mu,\sigma}(D)f(z)|^{1-\lambda} |\theta_{\mu,\sigma}(D)f(z)|^{\lambda}}{(1+\sigma^{-1}d((x,\omega),(z,\mu))^{2})^{N(1-\la)}(1+\sigma^{-1}d((x,\omega),(z,\mu))^{2})^{N\la}}\ud z\ud\mu\\
&\leq \big(M_{N,\sigma}^{*}(W_{\sigma}f)(x,\omega)\big)^{1-\lambda}
{\sigma}^{-n}\int_{\Sp}\frac{|\theta_{\mu,\sigma}(D)f(z)|^{\lambda}}{(1+\sigma^{-1}d((x,\omega),(z,\mu))^{2})^{N\lambda}}\ud z\ud\mu\\
&\leq \big(M_{N,\sigma}^{*}(W_{\sigma}f)(x,\omega)\big)^{1-\lambda}
\big(\Ma_{\lambda}(W_{\sigma}f)(x,\omega)\big)^{\lambda}.
\end{align*}
This proves the required statement if $M_{N,\sigma}^{*}(W_{\sigma}f)(x,\omega)<\infty$, and it remains to show that $M_{N,\sigma}^{*}(W_{\sigma}f)(x,\w)$ is finite whenever $\mathcal{M}_{\lambda}(W_{\sigma}f)(x,\omega)$ is finite.

To this end, first write 
\[
%\begin{align}\label{9-29-11}
\theta_{\nu,\sigma}(D)f(y)=\theta_{\nu,\sigma}(D)\wt{\Psi}(\sigma D)f(y)
=\int_{\Rn}\F^{-1}\theta_{\nu,\sigma}(y-z)\wt{\Psi}(\sigma D)f(z)\ud z
%\end{align}
\]
for $(y,\nu)\in\Sp$, where we recall from Remark \ref{rem:thetatilde} that $\wt{\Psi}\equiv1$ on $\supp(\Psi)$. Since $\wt{\Psi}\in \Sw(\Rn)$ and $f\in\Sw'(\Rn)$, there exists an $M\geq0$, dependent on $f$, such that
\[
%C_{\sigma,f}:=
\sup_{z\in\Rn}\lb z\rb^{-2M}|\wt{\Psi}(\sigma D)f(z)|<\infty.
\]
Hence Lemma \ref{lem:packetbounds} implies that
\begin{equation}\label{eq:decaydependent}
\begin{aligned}
%  \begin{align}\label{9-29-12}
&\lb y\rb^{-2M}|\theta_{\nu,\sigma}(D)f(y)|\\
&\lesssim\int_{\Rn} \lb y-z\rb^{2M}|\F^{-1}\theta_{\nu,\sigma}(y-z)|\ud z\sup_{z\in\Rn} \lb z\rb^{-2M}|\wt{\Psi}(\sigma D)f(z)|\\
 &\lesssim \sup_{z\in\Rn} \lb z\rb^{-2M}|\wt{\Psi}(\sigma D)f(z)|\lesssim 1,
% \end{align}
\end{aligned}
\end{equation}
where the implicit constants depend on $f$ and $\sigma$, but not on $(y,\nu)$.

Next, for $0<\veps<1$, denote 
\[
\sup_{(y,\nu)\in\Sp}\frac{\big|\theta_{\nu,{\sigma}}(D)f(y)\big|}{(1+\sigma^{-1}d((x,\omega)(y,\nu))^{2})^{N}(1+\veps\sigma^{-1} d((x,\omega),(y,\nu))^{2})^{M}}
\]
by $M_{N,\sigma,\veps}^{*}(W_{\sigma}f)(x,\omega)$.
Then, by \eqref{eq:equivmetric} and \eqref{eq:decaydependent},
\begin{equation}\label{eq:Mepsfinite}
\begin{aligned}
%\begin{align}\label{tilde{M}}
&M_{N,\sigma,\veps}^{*}(W_{\sigma}f)(x,\omega)\\
&\leq\sup_{(y,\nu)\in\Sp}(1+\veps\sigma^{-1} d((x,\omega),(y,\nu))^{2})^{-M}|\theta_{\nu,\sigma}(D)f(y)|\\
&\lesssim \sup_{(y,\nu)\in\Sp}\lb x-y\rb^{-2M}|\theta_{\nu,\sigma}(D)f(y)|\\
&\lesssim\lb x\rb^{2M}\sup_{(y,\nu)\in\Sp}\lb y\rb^{-2M}|\theta_{\nu,{\sigma}}(D)f(y)|<\infty,
%\end{align}
\end{aligned}
\end{equation}
for implicit constants dependent on $f$, $\sigma$ and $\veps$. 

We can now combine \eqref{eq:reprotheta} and Lemma \ref{lem:packetbounds} in the same way as before to write
\begin{align*}
&|\theta_{\nu,\sigma}(D)f(y)|\lesssim\sigma^{-n}
\int_{\Sp}(1+\sigma^{-1}d((y,\nu),(z,\mu))^{2})^{-N-M}|\theta_{\mu,{\sigma}}(D)f(z)|\ud z\ud\mu\\
&\leq \sigma^{-n}
\int_{\Sp}\frac{|\theta_{\mu,\sigma}(D)f(z)|}{(1+\sigma^{-1}d((y,\nu),(z,\mu))^{2})^{N}(1+\veps\sigma^{-1}d((y,\nu),(z,\mu))^{2})^{M}}\ud z\ud\mu,
\end{align*}
for an implicit constant independent of $(y,\nu)$ and $\veps$. The triangle inequality then yields
\begin{align*}
&(1+\sigma^{-1}d((x,\omega),(y,\nu))^{2})^{-N}(1+\veps\sigma^{-1}d((x,\omega),(y,\nu))^{2})^{-M}|\theta_{\nu,\sigma}(D)f(y)|\\
&\lesssim\sigma^{-n}\int_{\Sp}\frac{|\theta_{\mu,\sigma}(D)f(z)|}{(1+\sigma^{-1}d((x,\w),(z,\mu))^{2})^{N}(1+\veps\sigma^{-1}d((x,\w),(z,\mu))^{2})^{M}}\ud z\ud\mu\\
&\leq \big(M_{N,\sigma,\veps}^{*}(W_{\sigma}f)(x,\omega)\big)^{1-\lambda}
\sigma^{-n}\int_{\Sp}\frac{|\theta_{\mu,{\sigma}}(D)f(z)|^{\lambda}}{(1+\sigma^{-1}d((x,\omega),(z,\mu))^{2})^{N\lambda}}
 \ud z\ud\mu.
\end{align*}
This in turn implies that
\begin{align*}
M_{N,\sigma,\veps}^{*}(W_{\sigma}f)(x,\omega)\lesssim \big(M_{N,\sigma,\veps}^{*}(\W_{\sigma}f)(x,\omega)\big)^{1-\lambda}\big(\Ma_{\lambda}(W_{\sigma}f)(x,\omega)\big)^{\lambda},
\end{align*}
and combined with \eqref{eq:Mepsfinite} we now obtain
\[
M_{N,\sigma,\veps}^{*}(W_{\sigma}f)(x,\omega)\lesssim\Ma_{\lambda}(W_{\sigma}f)(x,\omega).
\]
Finally, since the implicit constant is independent of $\veps$, letting $\varepsilon$ tend to zero yields
\begin{align*}
M_{N,\sigma}^{*}(W_{\sigma}f)(x,\omega)\lesssim\Ma_{\lambda}(W_{\sigma}f)(x,\omega).
\end{align*}
This shows that $M_{N,\sigma}^{*}(W_{\sigma}f)(x,\omega)$ is indeed finite if $\Ma_{\lambda}(W_{\sigma}f)(x,\omega)$ is finite. %, which concludes the proof.
\end{proof}

\end{document}
